\documentclass[11pt,a4paper]{article}
\usepackage[utf8]{inputenc}
\usepackage[vietnamese]{babel}
\usepackage{amsmath,amssymb}
\usepackage{geometry}
\geometry{a4paper, total={180mm,260mm}, left=15mm, right=15mm, top=18mm, bottom=18mm}
\usepackage{tikz}
\usetikzlibrary{shapes.geometric, arrows.meta, positioning, calc}
\usepackage{tcolorbox}
\tcbuselibrary{skins,breakable}
\usepackage{xcolor}
\usepackage{multicol}
\usepackage{array}
\usepackage{fontawesome}
\usepackage{fancyhdr}

\pagestyle{fancy}
\fancyhf{}
\renewcommand{\headrulewidth}{0.5pt}
\renewcommand{\footrulewidth}{0.5pt}
\fancyhead[L]{\small\bfseries\color{teal}\faBook\ SỔ TAY KIẾN THỨC TOÁN 3}
\fancyhead[R]{\small\bfseries\color{teal}Lớp Toán Thầy Lĩnh -- loptoanthaylinh.vn}
\fancyfoot[C]{\small\color{gray}-- Sổ tay kiến thức: Bài 59. Các số có năm chữ số. Số 100 000 --}

\definecolor{primary}{HTML}{005F73}
\definecolor{secondary}{HTML}{0A9396}
\definecolor{accent}{HTML}{EE9B00}
\definecolor{danger}{HTML}{AE2012}
\definecolor{boxbg}{HTML}{F4F6F8}

\begin{document}

% --- BANNER TIÊU ĐỀ ---
\begin{tcolorbox}[
  enhanced,
  colback=white,
  colframe=primary,
  arc=3mm,
  boxrule=1.2pt,
  left=4mm, right=4mm, top=2.5mm, bottom=2.5mm,
  drop shadow=black!10!white
]
  \centering
  {\Large\bfseries\color{primary}\faBookmark\ BÀI 59: CÁC SỐ CÓ NĂM CHỮ SỐ. SỐ 100 000}\\[1.5mm]
  {\small\itshape\color{black!75}Cấu tạo số: Chục nghìn, Nghìn, Trăm, Chục, Đơn vị; đọc viết số và số 100 000 (Một trăm nghìn) trên 1 trang A4}
\end{tcolorbox}

\vspace{-2mm}


\noindent
\begin{minipage}[t]{0.49\textwidth}\vspace{0pt}
\begin{tcolorbox}[
  enhanced,
  colback=blue!4!white,
  colframe=blue!75!black,
  arc=2mm, boxrule=0.8pt,
  left=2mm, right=2mm, top=1.8mm, bottom=1.8mm,
  title={\footnotesize\bfseries\color{white}\faKey\ 1. Cấu tạo số có 5 chữ số}
]
{\footnotesize
\begin{itemize}\setlength{\itemsep}{1.5pt}\setlength{\leftskip}{-2.5mm}
  \item Số có 5 chữ số gồm các hàng tính từ phải sang trái:
  $$\mathbf{\text{Đơn vị}} \to \mathbf{\text{Chục}} \to \mathbf{\text{Trăm}} \to \mathbf{\text{Nghìn}} \to \mathbf{\text{Chục nghìn}}$$
  \item Phân tích cấu tạo số thập phân:
  $$\overline{abcde} = a \times 10\,000 + b \times 1000 + c \times 100 + d \times 10 + e$$
  \item \textbf{Số 100 000 (Một trăm nghìn / Mười vạn):}
  Là số nhỏ nhất có $6$ chữ số, là số liền sau của số $99\,999$:
  $$99\,999 + 1 = \mathbf{100\,000}$$
\end{itemize}
}
\end{tcolorbox}
\end{minipage}
\hfill
\begin{minipage}[t]{0.49\textwidth}\vspace{0pt}
\begin{tcolorbox}[
  enhanced,
  colback=teal!4!white,
  colframe=teal!75!black,
  arc=2mm, boxrule=0.8pt,
  left=2mm, right=2mm, top=1.8mm, bottom=1.8mm,
  title={\footnotesize\bfseries\color{white}\faSliders\ 2. Quy tắc đọc và viết số chuẩn}
]
{\footnotesize
\begin{itemize}\setlength{\itemsep}{1.5pt}\setlength{\leftskip}{-2.5mm}
  \item Tách số thành 2 lớp: \textbf{Lớp nghìn} (gồm hàng chục nghìn và nghìn) và \textbf{Lớp đơn vị} (trăm, chục, đơn vị).
  \item Đọc số lớp nghìn kèm chữ "nghìn", sau đó đọc tiếp số lớp đơn vị.
  \item Giữa lớp nghìn và lớp đơn vị để một khoảng trắng hẹp: ví dụ viết $35\,620$, không dùng dấu phẩy.
\end{itemize}
}
\end{tcolorbox}
\end{minipage}

\vspace{-2.5mm}
% --- KHỐI 3: SƠ ĐỒ TƯ DUY TIKZ ---
\noindent
\begin{tcolorbox}[
  enhanced,
  colback=orange!4!white,
  colframe=orange!80!black,
  arc=2mm, boxrule=0.8pt,
  left=2mm, right=2mm, top=1.5mm, bottom=1.5mm,
  title={\footnotesize\bfseries\color{white}\faSitemap\ SƠ ĐỒ TƯ DUY: HỆ THỐNG 5 HÀNG CỦA SỐ CÓ 5 CHỮ SỐ}
]
\centering
\begin{tikzpicture}[
  u/.style={draw=blue!80!black, fill=blue!10, rounded corners=2pt, font=\scriptsize\bfseries, minimum width=2.1cm, minimum height=0.65cm, align=center},
  arr/.style={-{Stealth[length=1.8mm]}, thick}
]
  \node[u, fill=purple!20, draw=purple!80!black] (h_cn) at (0, 0) {CHỤC NGHÌN\\$a \times 10\,000$};
  \node[u, fill=teal!20, draw=teal!80!black] (h_n) at (2.8, 0) {HÀNG NGHÌN\\$b \times 1000$};
  \node[u, fill=blue!20, draw=blue!80!black] (h_t) at (5.6, 0) {HÀNG TRĂM\\$c \times 100$};
  \node[u, fill=green!20, draw=green!70!black] (h_c) at (8.4, 0) {HÀNG CHỤC\\$d \times 10$};
  \node[u, fill=orange!20, draw=orange!80!black] (h_dv) at (11.2, 0) {HÀNG ĐƠN VỊ\\$e \times 1$};

  \draw[arr, purple!80!black] (h_cn) -- (h_n);
  \draw[arr, teal!80!black] (h_n) -- (h_t);
  \draw[arr, blue!80!black] (h_t) -- (h_c);
  \draw[arr, green!70!black] (h_c) -- (h_dv);
\end{tikzpicture}
\end{tcolorbox}

\vspace{-2.5mm}
% --- KHỐI 4: VÍ DỤ & MẸO TRÁNH BẪY ---
\noindent
\begin{minipage}[t]{0.49\textwidth}\vspace{0pt}
\begin{tcolorbox}[
  enhanced,
  colback=green!4!white,
  colframe=green!60!black,
  arc=2mm, boxrule=0.8pt,
  left=2mm, right=2mm, top=1.8mm, bottom=1.8mm,
  title={\footnotesize\bfseries\color{white}\faCheckCircle\ Ví dụ thực hành chuẩn mực}
]
{\footnotesize
\textbf{1. Đọc và viết số:}\\
-- Số gồm $6$ chục nghìn, $2$ nghìn, $5$ trăm, $0$ chục, $8$ đơn vị viết là: $\mathbf{62\,508}$.\\
Đọc là: \textit{Sáu mươi hai nghìn năm trăm linh tám}.\\
-- Số gồm $8$ chục nghìn, $0$ nghìn, $0$ trăm, $4$ chục, $0$ đơn vị viết là: $\mathbf{80\,040}$.\\
Đọc là: \textit{Tám mươi nghìn không trăm bốn mươi}.\\
\textbf{2. Số liền sau:}\\
Số liền sau của số $99\,999$ là $\mathbf{100\,000}$.
}
\end{tcolorbox}
\end{minipage}
\hfill
\begin{minipage}[t]{0.49\textwidth}\vspace{0pt}
\begin{tcolorbox}[
  enhanced,
  colback=red!4!white,
  colframe=red!70!black,
  arc=2mm, boxrule=0.8pt,
  left=2mm, right=2mm, top=1.8mm, bottom=1.8mm,
  title={\footnotesize\bfseries\color{white}\faExclamationTriangle\ Mẹo tránh bẫy trắc nghiệm}
]
{\footnotesize
\begin{itemize}\setlength{\itemsep}{1.5pt}\setlength{\leftskip}{-2.5mm}
  \item \textbf{Bẫy viết thiếu chữ số 0 khi đọc có chữ "linh" hoặc "không trăm":}\\
  "Hai mươi nghìn không trăm lẻ năm" phải viết là $\mathbf{20\,005}$ (đủ 5 chữ số), viết $2005$ là sai vì chỉ có 4 chữ số!
  \item Đếm kĩ số lượng chữ số trước khi kết luận.
\end{itemize}
}
\end{tcolorbox}
\end{minipage}


\end{document}
